%% kapa: exam=1 round=33 year=2022 period=2 form=A subject=accounting subject_ko=회계학 seq=33 number=73
%% source: Q-Net 감정평가사 기출문제 (cst00309) artlSeq=5209997 file=raw/1008/5209997/2022년도 제33회 감정평가사 1차시험 2교시(감정평가 관계법규, 회계학).pdf
%% conversion: kordoc 4.24.51 -> tools/gen1.py (auto); review=proofread-sonnet; flags=none
\begin{question}{73}
(주)감평은 20×1년 3월 제품 A(단위당 판매가격 ￦800) 1,000단위를 생산·판매하였다. 3월의 단위당 변동원가는 ￦500이고, 총고정원가는 ￦250,000이 발생하였다. 4월에는 광고비 ￦15,000을 추가 지출하면 ￦50,000의 매출이 증가할 것으로 기대하고 있다. 이를 실행할 경우 (주)감평의 4월 영업이익에 미치는 영향은? (단, 단위당 판매가격, 단위당 변동원가, 광고비를 제외한 총고정원가는 3월과 동일하다.)
\begin{choices}
  \item ￦3,750 감소
  \item ￦3,750 증가
  \item ￦15,000 감소
  \item ￦15,000 증가
  \item ￦35,000 증가
\end{choices}
\end{question}
\answer{2}
